%& -shell-escape
% !TeX root = thesis-text.tex
% !TeX encoding = UTF-8
% !TeX program = XeLaTeX

\documentclass{.local/lib/classes/ndvr-super-article-standalone}

\title[NDVR thesis-abstract]{%
    Проблема поиска нечетких дубликатов видео%
}

\begin{document}

    \section{Ранжирование результатов поиска}

    Обычно в~поисковой выдаче требуется группировать 
    одинаковые видео по~одним и~тем~же ключевым словам запроса.
    Но чаще всего поисковые машины возвращают список дубликатов.
    На рисунках~\ref{fig:video-reranking-mailru},
    \ref{fig:video-reranking-yandex},
    \ref{fig:video-reranking-vk} и
    \ref{fig:video-reranking-youtube} 
    показаны примеры неудачного ранжирования результатов
    для различных сервисов интернет-видео. 
    
    
    \begin{figure}[htp]
        \includegraphics[width=0.8\textwidth]{.local/img/video/apps/video-reranking-mailru.png}
        \caption{
            Неудачное ранжирование при~поиске «Faszination~Korallenriff~(2011)».
            Сервис Mail.ru.
        }
        \label{fig:video-reranking-mailru}
        \vspace{2em}
        \includegraphics[width=0.8\textwidth]{.local/img/video/apps/video-reranking-yandex.png}
        \caption{
            Неудачное ранжирование при~поиске «Faszination Korallenriff (2011)».
            Сервис Яндекс.
        }
        \label{fig:video-reranking-yandex}
    \end{figure}

    \begin{figure}[htp]
        \includegraphics[width=0.8\textwidth]{.local/img/video/apps/video-reranking-vk.png}
        \caption{
            Неудачное ранжирование при~поиске «Faszination~Korallenriff~(2011)».
            Сервис Вконтакте (\href{http://vk.com/}{vk.com}).
        }
        \label{fig:video-reranking-vk}
        \vspace{2em}
        \includegraphics[width=0.8\textwidth]{.local/img/video/apps/video-reranking-youtube.png}
        \caption{
            Ранжирование при~поиске «Faszination~Korallenriff~(2011)».\\
            Сервис YouTube.
        }
        \label{fig:video-reranking-youtube}
    \end{figure}

    \noindent
    Качество поисковой выдачи можно улучшить тремя способами.
    \begin{enumerate}
        \item   Оставить в поисковой выдачи только 
                самые свежие копии видео.
        \item   Удалить нечётки дубликаты 
                из~верхней части списка \cite{Wu:2007a}.
        \item   Сгруппировать оригинальные видео \cite{Huang:2010a}.
    \end{enumerate}

    В~первом и во втором случае новизна результатов
    в~рейтинговом списке значительно увеличится.
    В~то~же время, видео, которое интересовало пользователя,
    может быть расположено слишком далеко от~начала страницы.
    Похожие видео могут потребоваться пользователям как~дополнительный
    источник информации.
    Например, видео в~конкретном разрешении и~форматах,
    с~различными субтитрами, аудио и~пр.  \cite{Cherubini:2009}

    Во~третьем случае описанной проблемы нет.
    Новые результаты просто можно выделить тем или~иным образом.
    Для всех трёх вариантов требуется эффективный поиск
    нечётких дубликатов видео.

\end{document}